\section{The Authentication Tag as a Hash Output}
\label{sec:root-tag-hash-security}

Externally, a \TW{} encryption returns a ciphertext body and one authentication
tag. Internally, that tag is the final root-duplex output: leaf tags summarize
leaf chunks, the root aggregates those summaries with the context and root
ciphertext chunk, and the wrapper $\HTW$ returns the tag alone. This section
proves that the visible tag is a hash output for the full opening
$(K, U, A, M)$, not merely an AEAD verification value.

The hash notions are the standard ones for $\HTW$: collision,
preimage, second-preimage, and everywhere preimage resistance. The committing
statements are consequences of that hash security: a successful compact opening
must force distinct openings to agree on the same $\tauroot$-bit tag, the
root-collision event bounded below.

The proof uses three ingredients. Transcript parsing separates final root
inputs, except for the explicitly charged leaf-tag collision event.
\Cref{lem:root-candidate-seq} orders every relevant final-root input into a
chronological candidate sequence, and \Cref{lem:root-tag-hash} applies the
adaptive collision and preimage bound of \Cref{app:lemmas} to that sequence.
The random oracle theorems are short instantiations; the flat sponge lift gives
the public-permutation bounds; the \cite{BH22} committing notions reuse the
same root-collision bound.

\Cref{tab:proof-dependency-map} maps the final claims to the proof ingredients
and loss terms they use. It is only a dependency guide; the formal games and
resource accounting remain those of \Cref{sec:security-model}.

\begin{table}[t]
\centering
\footnotesize
\setlength{\tabcolsep}{3pt}
\begin{tabular}{@{}p{0.22\linewidth}p{0.70\linewidth}@{}}
\toprule
Claim & Dependency summary \\
\midrule
Coll, \Cref{thm:coll}
  & Terms: $\Lift_c$, $\LeafColl$, $\RootColl$. Ingredients: final-root input
    separation under no leaf collision; adaptive root candidate collision
    bound. \\
Pre, \Cref{thm:pre}
  & Terms: $\Lift_c$, $\LeafColl$, $\RootPre$, source-guess term.
    Ingredients: source/output separation under no leaf collision;
    sampled-source candidate hit bound. \\
ePre, \Cref{thm:epre}
  & Terms: $\Lift_c$, $\RootPre$. Ingredient: fixed-target root candidate hit
    bound. \\
CMTDs-4, \Cref{cor:cmtds4}
  & Terms: $\Lift_c$, $\RootColl$. Ingredients: \TW{} tidiness; common
    ciphertext body; opening-triple separation. \\
CMT/CMTD, \Cref{cor:committing-combined}
  & Terms: $\Lift_c$, $\RootColl$. Ingredients: \cite{BH22} hierarchy for
    D-notions; direct encryption-source root collision argument for CMTs. \\
mu-ind and AE, \Cref{thm:mu-ind,cor:ae}
  & Terms: $\Lift_c$, $\Lift_c(\Nprim)$, $\KeyGuess$, leaf tag, root tag.
    Ingredients: adjacent-user hybrid; hidden-key leaf collision and
    verification-first root guessing. \\
\textsf{INT-RUP}, \Cref{thm:int-rup}
  & Terms: $\Lift_c$, $\KeyGuess$, leaf tag, root tag. Ingredients:
    plaintext-computing work charged to transcript resources; hidden-key leaf
    collision and \textsf{INT-RUP} root guessing. \\
\bottomrule
\end{tabular}
\caption{Proof dependency map for the final \TW{} bounds. ``Leaf tag'' denotes
the linear hidden-key leaf term in the AE-facing bounds and $\LeafColl$ where
shown explicitly; ``root tag'' denotes the verification or root-preimage
guessing term for the relevant game.}
\label{tab:proof-dependency-map}
\end{table}

\subsection{Tag Hash and Committing Games}
\label{sec:root-tag-games}

The public-permutation games in this section use the primitive, validity, and
construction-oracle conventions of \Cref{sec:games-pp}. The random oracle
versions are obtained by the mechanical substitution of \Cref{sec:ro-model}.

\paragraph{Hash games.}
The hash games for the wrapper $\HTW$ (\Cref{def:htw}) have no
challenger-sampled key. In the $s$-way collision game, all inputs are
adversary-supplied: the adversary outputs $s$ inputs, and the challenger
returns $1$ iff they are pairwise distinct, in the \TW{} domain, and share
a root tag:
\begin{align*}
  &\pi \sample \Perm(\bits^{1600}), \quad
   (K_i, U_i, A_i, M_i)_{i=1}^s \sample \mathcal{A}^{\pi,\, \pi^{-1}}, \\
  &\text{return $0$ unless the inputs are pairwise distinct and in the \TW{} domain}, \\
  &\text{return } \bigl[\HTW(K_1, U_1, A_1, M_1) = \cdots =
     \HTW(K_s, U_s, A_s, M_s)\bigr],
\end{align*}
with advantage $\AdvColl_\TW(\mathcal{A}) = \Pr[\text{the game returns }
1]$; no nonce uniqueness is imposed. The \cite{RS04} standard and
everywhere second-preimage notions are used on fixed input slices; in
\Cref{cor:second-preimage-rs04} they are reduced explicitly to this
two-way collision game.
For fixed byte lengths $a,\mu$, the standard preimage game samples
$X^* \sample \mathcal{X}_{a,\mu}$ independently of the primitive, computes
$T^* = \HTW(X^*)$, gives $T^*$ to
$\mathcal{A}^{\pi,\, \pi^{-1}}$, and returns $1$ iff $\mathcal{A}$ outputs a
valid input $X$ in the \TW{} domain with $\HTW(X) = T^*$; the advantage is
$\AdvPreSlice{a}{\mu}_\TW(\mathcal{A})$. The
everywhere preimage game fixes a target tag
$T^* \in \bits^{\tauroot}$ before the primitive sampling and returns $1$
iff $\mathcal{A}(T^*)$ outputs an input $X$ in the domain with
$\HTW(X) = T^*$, the advantage $\AdvePre_\TW$ taken as the worst case over
query-independent targets.

\paragraph{Committing games.}
The \cite{BH22} committing games are instantiated as in
\Cref{sec:commit-defs}, with no challenger-sampled key. The $s$-way
CMTDs-$\ell$ (decryption) game has the adversary output a ciphertext
$C \concat T$ and $s$ tuples $(K_i, U_i, A_i, M_i)$ with pairwise distinct
$\WiC_\ell$-images, and returns $1$ iff
$\Dec_{K_i}[\pi](U_i, A_i, C \concat T) = M_i$ for every $i$, with the eager
length rejection justified by \Cref{alg:dec}
($\Dec_K(U, A, C \concat T)$ returns $\bot$ or a message of length $\|C\|$);
its advantage is
$\Advcmtds{\ell}_\TW(\mathcal{A})$. The CMTs-$\ell$ (encryption) game has
the adversary output the $s$ tuples alone and returns $1$ iff the
encryptions $\Enc_{K_i}[\pi](U_i, A_i, M_i)$ are all equal; its advantage is
$\Advcmts{\ell}_\TW(\mathcal{A})$.

\subsection{Separating Final Tag Inputs}

The tag value is read from a final root transcript input. The next two facts
rule out structural collisions at that input, except for the explicitly charged
possibility that two leaf tags collide before the root aggregates them.

\begin{proposition}[Opening Triple Separation]
\label{cor:triple-sep}
If two \TW{} root computations have distinct $(K, U, A)$ triples, then
their bridged final root transcript inputs under $\flatmap(\cdot)$ are
distinct, regardless of whether each computation is an encryption or
a decryption check, and whether or not they act on a common ciphertext.
\end{proposition}

\begin{proof}
Consider two root computations $i$ and $j$ with contexts
$(K_i, U_i, A_i) \neq (K_j, U_j, A_j)$, reaching final root transcripts
$R_i$ and $R_j$. Suppose for contradiction that
$\flatmap(R_i) = \flatmap(R_j)$. Since the
bridge $\flatmap(\cdot)$ is injective on well-formed transcript prefixes
(\Cref{lem:bridge-decode}), the native flattened inputs of the two
final-root prefixes are equal. Then \Cref{lem:transcript-inj}
recovers the full duplex-call sequence and reads off the pair
$(\sigma,\sigph)$ for each call. The seven $4$-bit phase suffixes are
pairwise distinct, so the \siginit, $\sigma_{\mathsf{ad}}^\pm$,
$\sigma_{\mathsf{msg}}^\pm$, and (when present) $\sigma_{\mathsf{agg}}^\pm$
calls in the recovered sequence are unambiguously identified. The
aggregation phase is present exactly when $n \geq 1$, and the two
recovered sequences are equal, so they either both carry an aggregation
phase or both omit it; either way the recovery of $(K, U, A)$ below uses only the
\siginit and AD calls, which precede any message or aggregation block.

Equality of the recovered sequences forces equality of the \siginit
call, whose $\sigma$ is $\prfxroot \concat K \concat U$ with
fixed-width $K$ and $U$, hence $K_i = K_j$ and $U_i = U_j$. It also
forces equality of the recovered AD subsequences. By the AD bullet of
\Cref{lem:transcript-inj} the associated data phase contributes a
$\sigadlast$-tagged block exactly when the associated data is nonempty,
so the recovered sequences agree on the presence of the AD phase. If
exactly one of $A_i, A_j$ were empty the recovered sequences would
differ in the presence of a $\sigadlast$ call, contradicting their
equality; hence either both are empty, giving $A_i = A_j = \varepsilon$,
or both are nonempty, in which case the phase-recovery sub-claim of
\Cref{lem:transcript-inj} inverts the $\rho$-byte-block partition map on
the AD byte string and equal AD $\sigma$-block sequences imply
$A_i = A_j$. In all cases $(K_i, U_i, A_i) = (K_j, U_j, A_j)$,
contradicting the hypothesis of distinct triples.
\end{proof}

Triple Separation handles distinct contexts; when inputs share a context
but differ in their messages, distinctness of the final root transcripts
needs the absence of a leaf tag collision, recorded next.

\begin{proposition}[Final-Root Input Separation]
\label{prop:final-root-sep}
Let $\badleaf^\star$ be the event that two distinct construction-generated
final leaf transcripts receive the same $\tauleaf$-bit $\RF$ projection.
On $\neg\badleaf^\star$, if two \TW{} construction computations reach
final root transcripts with $\flatmap(R) = \flatmap(R')$, then they share
the same context $(K, U, A)$ and the same full ciphertext $C$.
Equivalently, on $\neg\badleaf^\star$ computations with distinct
$(K, U, A, C)$ reach distinct bridged final root transcripts.
\end{proposition}

\begin{proof}
By bridge injectivity (\Cref{lem:bridge-decode}),
$\flatmap(R) = \flatmap(R')$ gives equal native flattened final-root
inputs. \Cref{lem:transcript-inj} then recovers the same root
initialization $\prfxroot \concat K \concat U$---hence the same $K$ and
$U$---the same associated data blocks---hence the same $A$---the same root
ciphertext chunk, and the same leaf count and ordered leaf tag sequence.
Under $\neg\badleaf^\star$, Leaf-Transcript Recovery
(\Cref{lem:leaf-recovery}) identifies the equal leaf tag sequences with
the same final leaf transcripts, hence the same leaf ciphertext bodies;
with the common root ciphertext chunk this gives the same full $C$.
\end{proof}

\subsection{Candidate Bound for Tag Inputs}

After separation, the probabilistic question is local: how likely is a collision
or target hit among the final root inputs whose tag values matter to a game?
The candidate machinery rests on one engine, stated here and proved in full in
\Cref{app:lem-candidate-collisions}: adaptively chosen distinct random oracle
inputs whose values are still unread when chosen behave, for collision and
target-hitting purposes, as fresh uniform draws.

\begin{lemma}[Adaptive Candidate Collision and Preimage]
\label{lem:adaptive-candidate}
Consider any interaction with a lazy random oracle $\ROflat : \bits^* \to
\bits^\infty$ that builds a sequence of distinct \emph{candidate} inputs
$Y_1,\dots,Y_m$. The process may make arbitrary non-candidate oracle
queries and may choose each candidate adaptively from the full prior
transcript and private state sampled before the candidate's unread oracle
value. Whenever a new candidate $Y_j$ is appended, two premises hold:
\emph{predictability}---both the decision to append and the string $Y_j$
are determined by the state preceding the append step---and
\emph{unrevealed value}---no oracle lookup before the append step has
requested a positive number of bits of $\ROflat(Y_j)$; zero-length lookups
at $Y_j$ are permitted. Fix a projection length $\tau$, and suppose the
sequence has at most $L$ candidates ($m \leq L$).
\begin{enumerate}
\item \emph{Collision.} For every $s \geq 2$, the
  probability $\Pr[\mathsf{coll}_{s,\tau}]$ that some $s$ distinct
  candidates have equal $\tau$-bit projections,
  \[
    \bigl\lfloor \ROflat(Y_{i_1}) \bigr\rfloor_\tau = \cdots =
    \bigl\lfloor \ROflat(Y_{i_s}) \bigr\rfloor_\tau
    \quad \text{for some } 1 \leq i_1 < \cdots < i_s \leq m,
  \]
  satisfies
  $\Pr[\mathsf{coll}_{s,\tau}] \leq \binom{L}{s} \cdot 2^{-\tau(s-1)}$.
\item \emph{Preimage.} For every target $t \in \bits^\tau$ chosen
  independently of $\ROflat$, the probability
  $\Pr[\mathsf{pre}_{\tau,t}]$ that some candidate's $\tau$-bit projection
  equals $t$, i.e.\ $\bigl\lfloor \ROflat(Y_j) \bigr\rfloor_\tau = t$ for
  some $1 \leq j \leq m$, satisfies
  $\Pr[\mathsf{pre}_{\tau,t}] \leq L \cdot 2^{-\tau}$.
\end{enumerate}
The preimage case has a \emph{second-preimage} corollary: for a
designated candidate position $j_0$, the probability
$\Pr[\mathsf{sec}_{\tau,j_0}]$ that some candidate $j \neq j_0$ has
$\bigl\lfloor \ROflat(Y_j) \bigr\rfloor_\tau =
 \bigl\lfloor \ROflat(Y_{j_0}) \bigr\rfloor_\tau$, i.e.\ collides with the
designated candidate on its $\tau$-bit projection, satisfies
$\Pr[\mathsf{sec}_{\tau,j_0}] \leq (L-1) \cdot 2^{-\tau}$.
\end{lemma}

\begin{proof}[Proof sketch]
The premises license lazy sampling at the append step: predictability
fixes the candidate string before its value is drawn, and the
unrevealed-value premise states that no bit of that value has been
exposed, so $\ROflat(Y_j)$ is uniform and independent of the entire prior
transcript. Conditioning position by position, each fixed position set
then contributes at most $2^{-\tau(s-1)}$ (each fixed position at most
$2^{-\tau}$), and a union bound over the $\binom{L}{s}$ position sets
(the $L$ positions) gives the bounds. The full argument, including the
second-preimage corollary, appears in
\Cref{app:lem-candidate-collisions}.
\end{proof}

The remaining work is structural: exhibiting, in each game, a candidate
sequence that satisfies these premises and contains the final root
transcripts the win event constrains.

\begin{lemma}[Root Tag Candidate Sequence]
\label{lem:root-candidate-seq}
Consider one of the hash or committing random oracle games of
\Cref{sec:ro-model}. The adversary $\mathcal{B}$ makes direct
$\ROflat$ queries and has no construction oracle. The only
construction-internal positive-length lookups at root tag inputs are the
root tag samplings of $h \geq 1$ deterministic challenger encryption or
decryption checks at game-fixed points. Apart from the standard preimage
game's designated sampled source check, all check inputs and keys are
adversary supplied. In that preimage game the sampled source input,
including its key, is private challenger state sampled before the source
tag is revealed, and the source computation counts as one of the $h$
checks. Let the $i$-th check reach a final root transcript $R_i$ whose
root tag is read by a single $\ROflat$ lookup at $\flatmap(R_i)$. Call an
$\ROflat$ input a \emph{root tag input} if it decodes under
$\KeyedTWOut$ to a root tag output-point class---$\Rmsg^+$ (the final
root message block, which carries the root tag only on the no-leaf path)
or $\Rtag$ (the aggregated root tag); multi-chunk computations also reach
$\Rmsg^+$ with zero requested output before aggregation, so this predicate
over-approximates the root tag candidate set, which only loosens the
count. Every execution
then admits a chronological candidate sequence in which each root tag
input is appended exactly once, immediately before the first event that is
either a direct $\ROflat$ query at it or a positive-length $\ROflat$
lookup reading it, and inputs never reached by such an event are never
appended. This sequence is such that
\begin{enumerate}
\item the candidates are pairwise distinct, and if the challenger reaches
  its checks, each $\flatmap(R_i)$ is in the sequence no later than the
  sampling of the $i$-th root tag answer;
\item it satisfies the predictability and unrevealed-value premises of
  \Cref{lem:adaptive-candidate}; and
\item its length is at most $\qRO(\mathcal{B}) + h$.
\end{enumerate}
\end{lemma}

\begin{proof}
A direct query to a keyed transcript is one of $\mathcal{B}$'s visible
$\ROflat$ queries counted in $\qRO(\mathcal{B})$, not a separate
construction query. The only construction computations in the game are
the $h$ deterministic challenger checks named in the statement.

\emph{Distinctness and coverage.} Each root tag input is appended at its
first triggering event and never again, so the candidates are pairwise
distinct by construction, and the rule is well defined however direct
queries and challenger checks interleave. In particular, in the standard
preimage game the source check's root tag sampling may be the trigger for
$\flatmap(R_1)$; a later direct query at the same address is not a first
trigger and appends nothing. Each $\flatmap(R_i)$ is appended no later
than the sampling of the $i$-th root tag answer, since that sampling is a
positive-length lookup at $\flatmap(R_i)$, giving~(1).

\emph{Predictability.} The decoded root-tag-input predicate depends on
the input string alone through the exact decoder of
\Cref{def:keyed-twout}, which ignores the requested output length; by
Output-Point Separation (\Cref{cor:output-sep}) each input decodes to at
most one output-point class. The append rule also uses the already fixed
next event: a direct query triggers at the address, while a construction
lookup triggers only when it requests positive output length. The string
itself is fixed in the state preceding the trigger: a direct query's
input is chosen by $\mathcal{B}$ before its lookup, and a check's final
root transcript is determined by the check's input and the prior oracle
answers---including any private challenger state sampled before the
candidate's unread value, which \Cref{lem:adaptive-candidate} expressly
allows.

\emph{Unrevealed value.} The trigger is the first direct query or
positive-length lookup at the address, so no earlier direct query at it
exists, and it suffices that no earlier construction lookup has read a
positive number of its bits. An earlier positive-length construction
lookup is either a stream output point, a leaf-tag point $\Ltag(j)$, or a
root-tag sampling. The first two are different output-point classes and
hence, by \Cref{cor:output-sep}, have different addresses from root tag
inputs; the last would itself have been the first trigger at this
address. Zero-requested-output lookups, including the $\Rmsg^+$ lookup a
multi-chunk computation makes before aggregation, are permitted by the
unrevealed-value premise of \Cref{lem:adaptive-candidate}, including
when such a zero-length lookup lands at another check's future no-leaf
final-root address. Together with predictability this gives~(2).

\emph{Length.} Appends are triggered only by direct queries---at most
$\qRO(\mathcal{B})$ of them---or by the $h$ checks' root tag samplings;
zero-length construction lookups never trigger an append. Hence the
sequence has
length at most $\qRO(\mathcal{B}) + h$, giving~(3).
\end{proof}

\begin{lemma}[Root Tag Hash and Target-Hit Bound]
\label{lem:root-tag-hash}
Consider one of the hash or committing \TW{} random oracle games
(\Cref{sec:ro-model}) for an adversary $\mathcal{B}$, satisfying the
hypotheses of \Cref{lem:root-candidate-seq}, with win event
$\mathsf{W}$. By the Root
Tag Candidate Sequence lemma
(\Cref{lem:root-candidate-seq}) the root tag inputs form a chronological
candidate sequence meeting the premises of \Cref{lem:adaptive-candidate} with
$\tau = \tauroot$. Writing
$\qRO = \qRO(\mathcal{B})$:
\begin{enumerate}
\item \emph{Root-collision case.} If $\mathsf{W}$ forces $s$ checks to
  reach pairwise distinct bridged final root transcripts carrying a
  common root tag, then, with $h = s$ checks,
  $\Pr[\mathsf{W}] \leq \RootColl_{\tauroot}(\qRO, s)$;
\item \emph{Fixed-target case.} If $\mathsf{W}$ forces a single check
  ($h = 1$) to carry a root tag equal to a target fixed independently of
  $\ROflat$, then
  $\Pr[\mathsf{W}] \leq \RootPre_{\tauroot}(\qRO)$;
\item \emph{Sampled-source case.} If the game first runs a designated
  source check, reveals its root tag, and $\mathsf{W}$ forces one later
  check to reach a distinct bridged final root transcript carrying that
  designated tag, then, with $h = 2$ checks,
  $\Pr[\mathsf{W}] \leq \RootPre_{\tauroot}(\qRO)$.
\end{enumerate}
\end{lemma}

\begin{proof}
In each case the candidate sequence has the stated length and meets the
premises of \Cref{lem:adaptive-candidate}.
\emph{(1)} The $s$ distinct candidates carrying a common
$\tauroot$-projection are an $s$-way collision; the collision case with
$L = \qRO + s$ gives
$\binom{\qRO + s}{s} \cdot 2^{-\tauroot(s-1)}
= \RootColl_{\tauroot}(\qRO, s)$.
\emph{(2)} The candidate hitting the fixed target is a preimage; the
preimage case with $L = \qRO + 1$ gives
$(\qRO + 1) \cdot 2^{-\tauroot} = \RootPre_{\tauroot}(\qRO)$.
\emph{(3)} Let $j_{\mathsf{src}}$ be the position of the designated
source check. The later check is required to be a distinct bridged final
root transcript, so it is a candidate position
$j \neq j_{\mathsf{src}}$ with the same
$\tauroot$-bit projection. This is an application of the second-preimage
corollary, not the fixed-target case, because the designated source tag is
sampled from $\ROflat$ rather than fixed independently of it.
The candidate sequence has length at most $L = \qRO + 2$, and the
second-preimage corollary of \Cref{lem:adaptive-candidate} gives
$(L-1)2^{-\tauroot} = (\qRO+1)2^{-\tauroot}
= \RootPre_{\tauroot}(\qRO)$.
\end{proof}

\subsection{Random Oracle Tag-Hash Bounds}

Each random oracle bound now matches a hash game to a case of
\Cref{lem:root-tag-hash}. Collision and fixed-slice preimage resistance also
carry a leaf-collision term, because distinct openings may have different
ciphertext bodies before reaching the root.

\begin{theorem}[Random Oracle Collision Resistance of $\HTW$]
\label{thm:coll-ro}
For every $s \geq 2$ and every $s$-way collision
adversary $\mathcal{B}$ in the \TW{} random oracle model,
\[
  \AdvColl_{\RO}(\mathcal{B})
  \;\leq\;
  \LeafColl_{\tauleaf}(\qRO(\mathcal{B})+\nleaf(\mathcal{B}))
  + \RootColl_{\tauroot}(\qRO(\mathcal{B}), s).
\]
\end{theorem}

\begin{proof}
Run the random oracle collision game. Let $\badleaf^\star$ be the event
that two distinct construction-generated final leaf transcripts receive
the same $\tauleaf$-bit $\RF$ projection; by the general count
of \Cref{lem:leaf-collision},
$\Pr[\badleaf^\star] \leq \LeafColl_{\tauleaf}(\qRO(\mathcal{B})+\nleaf(\mathcal{B}))$,
so it remains to bound the win under $\neg\badleaf^\star$. After rejecting
unless the $s$ output inputs are pairwise distinct and in the \TW{}
domain, the challenger performs $s$ deterministic encryptions
$\Enc_{K_i}^{\RO}(U_i, A_i, M_i)$, reaching $R_1,\dots,R_s$, all carrying a
common root tag $T$ on a win. Pairwise distinct tuples $(K_i, U_i, A_i, M_i)$
have pairwise distinct $(K_i, U_i, A_i, C_i)$: two tuples differing in their
context already differ in $(K, U, A)$, while two sharing a context but
differing in the message yield distinct ciphertexts, since for a fixed
context encryption maps the message to the ciphertext bijectively. So on
$\neg\badleaf^\star$ Final-Root Input Separation
(\Cref{prop:final-root-sep}) makes $\flatmap(R_1),\dots,\flatmap(R_s)$
pairwise distinct. The Root Tag Hash and Target-Hit Bound (\Cref{lem:root-tag-hash},
item~1) bounds the win under $\neg\badleaf^\star$ by
$\RootColl_{\tauroot}(\qRO(\mathcal{B}),s)$; adding $\Pr[\badleaf^\star]$
gives the theorem.
\end{proof}

\begin{theorem}[Random Oracle Preimage Resistance of $\HTW$]
\label{thm:pre-ro}
Fix byte lengths $0 \leq a \leq A_{\max}$ and
$0 \leq \mu \leq M_{\max}$, and let
$m_{a,\mu}=k+\nu+8a+8\mu$. For every Pre adversary
$\mathcal{B}$ against $\HTW$ on the fixed input slice
$\mathcal{X}_{a,\mu}$ in the \TW{} random oracle model
(\Cref{sec:ro-model}),
\[
\begin{aligned}
  \AdvPreSlice{a}{\mu}_\RO(\mathcal{B})
  \;\leq\;&
  \LeafColl_{\tauleaf}(\qRO(\mathcal{B})+\nleaf(\mathcal{B})) \\
  &+ \RootPre_{\tauroot}(\qRO(\mathcal{B}))
  + 2^{\tauroot-m_{a,\mu}} .
\end{aligned}
\]
\end{theorem}

\begin{proof}
Run the random oracle Pre game. The challenger samples
$X^* \sample \mathcal{X}_{a,\mu}$, writes
$X^*=(K^*,U^*,A^*,M^*)$, computes $T^*=\HTW(X^*)$, gives $T^*$ to
$\mathcal{B}$, and later checks the output $X=(K,U,A,M)$. Let
$\mathsf{same}$ be the event $X=X^*$. For a
fixed random oracle table and fixed adversary coins, the adversary's output
is a function of the $\tauroot$-bit challenge tag $T^*$, so over the
uniform source $X^*$ it can equal $X^*$ with probability at most
$2^{\tauroot}/2^{m_{a,\mu}}=2^{\tauroot-m_{a,\mu}}$. Hence
$\Pr[\mathsf{same}] \leq 2^{\tauroot-m_{a,\mu}}$.

It remains to bound winning executions with $X \neq X^*$. Let
$\badleaf^\star$ be the event that two distinct construction-generated
final leaf transcripts in the source and output checks receive the same
$\tauleaf$-bit $\RF$ projection. By the general count of
\Cref{lem:leaf-collision}---which charges every direct query and every
construction-generated final leaf transcript regardless of key
provenance, covering the challenger-sampled source key---
\[
  \Pr[\badleaf^\star]
  \leq
  \LeafColl_{\tauleaf}(\qRO(\mathcal{B})+\nleaf(\mathcal{B})).
\]
On $\neg\badleaf^\star$, the distinct tuples $X^*$ and $X$ have distinct
final root transcripts: distinct contexts are separated by
\Cref{cor:triple-sep}, while equal contexts and distinct messages produce
distinct ciphertexts, and then \Cref{prop:final-root-sep} applies. A win
with $X \neq X^*$ therefore forces the later output check to hit the
designated source tag $T^*$ at a distinct bridged final root transcript.
The sampled-source case, item~3, of \Cref{lem:root-tag-hash} bounds this
probability by $\RootPre_{\tauroot}(\qRO(\mathcal{B}))$. Adding the
source-guess and leaf-collision events gives the theorem.
\end{proof}

\begin{theorem}[Random Oracle Everywhere Preimage Resistance of $\HTW$]
\label{thm:epre-ro}
For every query-independent target tag $T^* \in \bits^{\tauroot}$ and
every ePre adversary $\mathcal{B}$ in the \TW{}
random oracle model (\Cref{sec:ro-model}),
\[
  \AdvePre_\RO(\mathcal{B})
  \;\leq\; \RootPre_{\tauroot}(\qRO(\mathcal{B}))
  = \frac{\qRO(\mathcal{B}) + 1}{2^{\tauroot}}.
\]
\end{theorem}

\begin{proof}
Run the random oracle preimage game for the query-independent target
$T^*$, which is independent of $\ROflat$. After rejecting unless the
output input $X = (K, U, A, M)$ is in the \TW{} domain, the challenger
performs one encryption $\Enc_{K}^\RO(U, A, M)$, reaching a final root
transcript $R$. A win means $\HTW(X) = T^*$, i.e.\ the root tag
$\lfloor \ROflat(\flatmap(R)) \rfloor_{\tauroot}$ at $R$ equals the fixed
target $T^*$. The fixed-target case, item~2, of
\Cref{lem:root-tag-hash} gives
$\AdvePre_\RO(\mathcal{B}) \leq \RootPre_{\tauroot}(\qRO(\mathcal{B}))$.
\end{proof}

\subsection{Public Permutation Tag-Hash Bounds}
\label{sec:root-tag-pp}

These are the public-permutation hash claims for the authentication tag. The
flat sponge lift transfers the preceding random oracle bounds to the public
permutation model. For these one-world hash games, the lift contributes
$\Lift_c(\Ntot)$, and the derived random oracle adversary has
$\qRO(\mathcal{B}) \leq \Nprim$; challenger construction
computations---including the preimage game's source encryption---are
construction-internal work counted in $N$.

\begin{theorem}[Collision Resistance of $\HTW$]
\label{thm:coll}
For every $s \geq 2$ and every adversary $\mathcal{A}$ against the $s$-way
multi-collision resistance of $\HTW$ in the sense of \cite{BH22}, with the
resources of \Cref{sec:resources},
\[
  \AdvColl_\TW(\mathcal{A})
  \;\leq\;
  \Lift_c(\Ntot)
  + \LeafColl_{\tauleaf}(\Nprim+\Nleaf)
  + \RootColl_{\tauroot}(\Nprim, s).
\]
\end{theorem}

\begin{proof}
By the Lift Application corollary (\Cref{cor:lift-application}) for the
collision game, the derived random oracle adversary $\mathcal{B}$ has
$\qRO(\mathcal{B}) \leq \Nprim$ and $\nleaf(\mathcal{B}) \leq \Nleaf$. The
random oracle bound of \Cref{thm:coll-ro} is nondecreasing in both
$\qRO(\mathcal{B})$ and $\nleaf(\mathcal{B})$, so the monotone form of
\Cref{cor:lift-application} gives the stated bound. The $s$ colliding
inputs may differ in their messages, hence in their ciphertext bodies; the
$\LeafColl$ term accounts for the only way two distinct inputs sharing a
context $(K, U, A)$ can reach the same final root transcript---a collision
among the leaf tags it absorbs---while distinct contexts already reach
distinct final root transcripts by Opening Triple Separation
(\Cref{cor:triple-sep}).
\end{proof}

\begin{corollary}[Two-Way Collision Resistance of $\HTW$]
\label{cor:coll-rs04}
Specializing \Cref{thm:coll} to $s = 2$ gives the classical two-way
collision resistance of \cite{RS04}: for every collision adversary
$\mathcal{A}$ against $\HTW$ with the resources of \Cref{sec:resources},
\[
  \AdvColl_\TW(\mathcal{A})
  \;\leq\;
  \Lift_c(\Ntot)
  + \LeafColl_{\tauleaf}(\Nprim+\Nleaf)
  + \RootColl_{\tauroot}(\Nprim, 2),
\]
with $\RootColl_{\tauroot}(\Nprim, 2) = \binom{\Nprim + 2}{2} / 2^{\tauroot}$.
\end{corollary}

\begin{proof}
The $s = 2$ case of \Cref{thm:coll}; \cite{BH22} note that $s = 2$
recovers the classical notion of collision resistance, which under the
keyed-by-the-permutation convention of \Cref{sec:commit-defs} is the
Coll notion of \cite{RS04}.
\end{proof}

\begin{corollary}[Second-Preimage Resistance of $\HTW$]
\label{cor:second-preimage-rs04}
Fix byte lengths $0 \leq a \leq A_{\max}$ and
$0 \leq \mu \leq M_{\max}$. For every adversary against $\HTW$, with the
resources of \Cref{sec:resources}, in either the standard
second-preimage (Sec) or everywhere second-preimage (eSec) sense of
\cite{RS04} on the fixed input slice $\mathcal{X}_{a,\mu}$, the
corresponding advantage is bounded, independently of $(a, \mu)$, by the
two-way collision bound of
\Cref{cor:coll-rs04}:
\[
  \Lift_c(\Ntot)
  + \LeafColl_{\tauleaf}(\Nprim+\Nleaf)
  + \RootColl_{\tauroot}(\Nprim, 2).
\]
\end{corollary}

\begin{proof}
We use \cite{RS04} to identify the standard and everywhere second-preimage
notions with the fixed-slice notions
$\mathrm{Sec}[m_{a,\mu}]$ and $\mathrm{eSec}[m_{a,\mu}]$ under the
canonical encoding of $\mathcal{X}_{a,\mu}$; the quantitative reduction
is explicit.

For a standard Sec adversary $\mathcal{A}$, define a two-way collision
adversary $\mathcal{C}$ as follows. It samples
$X_0 \sample \mathcal{X}_{a,\mu}$, runs
$\mathcal{A}^{\pi,\pi^{-1}}(X_0)$, receives $X_1$, and outputs
$(X_0,X_1)$. Whenever $\mathcal{A}$ wins, $X_1$ is a valid input,
$X_1 \neq X_0$, and $\HTW(X_1)=\HTW(X_0)$, so $\mathcal{C}$ wins the
two-way collision game with exactly the same probability.

For eSec, fix a source $X_0 \in \mathcal{X}_{a,\mu}$ and define
$\mathcal{C}_{X_0}$ to run the eSec adversary on that source and output
$(X_0,X_1)$. The same pointwise implication holds for every $X_0$; taking
the maximum over sources gives the eSec bound.

The induced collision adversary makes exactly the direct primitive queries
made by $\mathcal{A}$. Its two deterministic challenger checks are the
source and candidate checks, counted exactly as in the $s=2$ collision
game. Applying \Cref{cor:coll-rs04} to the induced collision adversary
gives the displayed bound.
\end{proof}

\begin{theorem}[Preimage Resistance of $\HTW$]
\label{thm:pre}
Fix byte lengths $0 \leq a \leq A_{\max}$ and
$0 \leq \mu \leq M_{\max}$, and let
$m_{a,\mu}=k+\nu+8a+8\mu$. For every adversary $\mathcal{A}$ against
$\HTW$, with the resources of \Cref{sec:resources}, in the standard
preimage (Pre) sense of \cite{RS04} on the fixed input slice
$\mathcal{X}_{a,\mu}$,
\[
\begin{aligned}
  \AdvPreSlice{a}{\mu}_\TW(\mathcal{A})
  \;\leq\;&
  \Lift_c(\Ntot)
  + \LeafColl_{\tauleaf}(\Nprim+\Nleaf) \\
  &+ \RootPre_{\tauroot}(\Nprim)
  + 2^{\tauroot - m_{a,\mu}}.
\end{aligned}
\]
\end{theorem}

\begin{proof}
By the Lift Application corollary (\Cref{cor:lift-application}) for the
preimage game, the derived random oracle adversary $\mathcal{B}$ has
$\qRO(\mathcal{B}) \leq \Nprim$ and
$\nleaf(\mathcal{B}) \leq \Nleaf$; per \Cref{sec:resources}, these caps
cover both challenger encryptions' final leaf transcripts, the source
encryption's included; the encryptions' duplex calls are counted in $N$
and enter through $\Lift_c(\Ntot)$. The random oracle bound of
\Cref{thm:pre-ro} is nondecreasing in both resources, so the monotone form
of \Cref{cor:lift-application} gives the stated bound.
\end{proof}

\begin{theorem}[Everywhere Preimage Resistance of $\HTW$]
\label{thm:epre}
For every target tag $T^* \in \bits^{\tauroot}$ fixed in advance and every
ePre adversary $\mathcal{A}$ against $\HTW$ with the resources of
\Cref{sec:resources},
\[
  \AdvePre_\TW(\mathcal{A})
  \;\leq\;
  \Lift_c(\Ntot) + \RootPre_{\tauroot}(\Nprim)
  \;=\;
  \Lift_c(\Ntot) + \frac{\Nprim + 1}{2^{\tauroot}}.
\]
\end{theorem}

\begin{proof}
By the Lift Application corollary (\Cref{cor:lift-application}) for the
everywhere-preimage game---whose target $T^*$ is fixed before the
primitive
sampling---the derived random oracle adversary $\mathcal{B}$ has
$\qRO(\mathcal{B}) \leq \Nprim$, and the random oracle bound of
\Cref{thm:epre-ro} is nondecreasing in $\qRO(\mathcal{B})$, so the monotone
form of \Cref{cor:lift-application} gives the stated bound. Compared with the
$s$-way root-collision term $\RootColl_{\tauroot}(\cdot, s)$ of
\Cref{sec:loss-terms}, the single fixed target replaces the exponent
$\tauroot(s-1)$ by $\tauroot$.
\end{proof}

\begin{corollary}[Tag Hash Security and Compact Commitment]
\label{cor:tag-commitment}
The \TW{} authentication tag, represented by the wrapper $\HTW$, is a secure
hash of the full encryption context $(K, U, A, M)$ in the sense of
\cite{RS04}: collision-resistant
(\Cref{thm:coll}), preimage-resistant on each fixed input slice
(\Cref{thm:pre}), second-preimage-resistant
(\Cref{cor:second-preimage-rs04}), and everywhere preimage-resistant
(\Cref{thm:epre}). Under the implemented
parameters and work cap $\Ntot \leq 2^{63}$, these advantages are at most
$2^{-128}$ (\Cref{cor:concrete-security}). These hash bounds make the
single $\tauroot$-bit tag a binding compact commitment to $(K, U, A, M)$, with
standard and everywhere preimage resistance, native to the mode. The commitment
is to an opening supplied for the tag: it includes the key and does not make the
tag an independently verifiable public digest of $(U,A,M)$.
\end{corollary}

\begin{proof}
Collision resistance is \Cref{thm:coll}; standard preimage resistance on
fixed input slices is \Cref{thm:pre}; second-preimage resistance follows from
collision resistance by \Cref{cor:second-preimage-rs04}; and everywhere
preimage resistance is \Cref{thm:epre}. These bounds quantify the
adversary's success at the same $\tauroot$-bit projection of the final root
transcript, and \Cref{cor:concrete-security} evaluates them at the
implemented parameters. Binding is the compact-commitment property supplied
by collision and second-preimage resistance. Standard preimage resistance is
the property that a tag produced from a uniformly sampled fixed-length
opening does not reveal an efficiently findable opening; everywhere preimage
resistance is the corresponding fixed-target property for tags chosen
independently of the primitive.
\end{proof}

\subsection{Committing Consequences of Tag Hash Security}

The committing consequences use the same tag-collision analysis under the
opening constraints of the \cite{BH22} games. CMTDs-4 is slightly tighter than
general collision resistance: all openings share one ciphertext body, so
distinct openings separate through their opening triples and no leaf-collision
term is needed.

\begin{theorem}[Random Oracle CMTDs-4]
\label{thm:cmtds4-ro}
For every $s \geq 2$ and every $s$-way CMTDs-4
adversary $\mathcal{B}$ in the \TW{} random oracle model
(\Cref{sec:ro-model}),
\[
  \Advcmtds{4}_\RO(\mathcal{B})
  \;\leq\; \RootColl_{\tauroot}(\qRO(\mathcal{B}), s)
  = \frac{\binom{\qRO(\mathcal{B}) + s}{s}}{2^{\tauroot (s - 1)}}.
\]
\end{theorem}

\begin{proof}
Run the random oracle CMTDs-4 game. After rejecting on any failed syntax,
message-length, or pairwise-full-input check, the challenger performs $s$
deterministic decryption checks $\Dec_{K_i}^\RO(U_i, A_i, C \concat T)$ on
the one common ciphertext $C \concat T$, reaching $R_1,\dots,R_s$. On a winning
execution the tuples $(K_i, U_i, A_i, M_i)$ are pairwise distinct, and since
all checks decrypt the common ciphertext, determinism of $\Dec$ makes the
triples $(K_i, U_i, A_i)$ pairwise distinct---equal triples would return
equal messages, hence equal tuples. By Opening Triple Separation
(\Cref{cor:triple-sep}) the candidates $\flatmap(R_1),\dots,\flatmap(R_s)$
are pairwise distinct (leaf tag collisions cannot merge them), and every
accepting check carries the common root tag $T$. The Root Tag Hash and Target-Hit Bound
(\Cref{lem:root-tag-hash}, item~1) gives
$\Advcmtds{4}_\RO(\mathcal{B}) \leq \RootColl_{\tauroot}(\qRO(\mathcal{B}), s)$.
\end{proof}

\begin{corollary}[CMTDs-4 Committing Security]
\label{cor:cmtds4}
For every $s \geq 2$ and every $s$-way CMTDs-4
adversary $\mathcal{A}$ against \TW{} with the resources of
\Cref{sec:resources},
\[
  \Advcmtds{4}_\TW(\mathcal{A})
  \;\leq\;
  \Lift_c(\Ntot) + \RootColl_{\tauroot}(\Nprim, s)
  \;=\;
  \Lift_c(\Ntot) + \frac{\binom{\Nprim + s}{s}}{2^{\tauroot (s - 1)}}.
\]
\end{corollary}

\begin{proof}
A CMTDs-4 winner gives one ciphertext $C \concat T$ and $s$ distinct inputs
$(K_i, U_i, A_i, M_i)$ such that
$\Dec_{K_i}(U_i, A_i, C \concat T) = M_i$ for every $i$. By \TW{} tidiness
(\Cref{lem:tw-tidiness}), each opened input re-encrypts to the same ciphertext:
$\Enc_{K_i}(U_i, A_i, M_i) = C \concat T$. In particular the opened inputs share
the root tag $T$. For the displayed no-$\LeafColl$ bound, use the dedicated
root-collision derivation of \Cref{thm:cmtds4-ro}, lifted by
\Cref{cor:lift-application}: in this all-bodies-equal case, the single shared
body $C$ forces the triples $(K_i, U_i, A_i)$ pairwise distinct; by Opening
Triple Separation (\Cref{cor:triple-sep}), a leaf tag collision cannot merge
two openings and no $\LeafColl$ term arises.
\end{proof}

\paragraph{Remaining committing notions.}
The hierarchy and $s$-way extension of the \cite{BH22} committing notions
are given in \cite[\S3, Figs.~4--5, and App.~A]{BH22}. We use those
relations for the D-notions and a direct source-game argument for the
CMTs-$\ell$ notions to preserve the resource split of
\Cref{sec:resources}.

\begin{corollary}[Committing-Notion Root-Collision Bound]
\label{cor:committing-combined}
For every $s \geq 2$, every
\[
  X \in \bigl\{\mathsf{CMTDs\text{-}1},\, \mathsf{CMTDs\text{-}3},\,
              \mathsf{CMTDs\text{-}4},\, \mathsf{CMTs\text{-}1},\,
              \mathsf{CMTs\text{-}3},\, \mathsf{CMTs\text{-}4}\bigr\},
\]
and every $s$-way $X$-adversary $\mathcal{A}$ against \TW{} with the
resources of \Cref{sec:resources},
\[
  \Adv^X_\TW(\mathcal{A})
  \;\leq\;
  \Lift_c(\Ntot) + \RootColl_{\tauroot}(\Nprim, s)
  \;=\;
  \Lift_c(\Ntot) + \frac{\binom{\Nprim + s}{s}}{2^{\tauroot(s-1)}}.
\]
\end{corollary}

\begin{proof}
CMTDs-4 is bounded directly by \Cref{cor:cmtds4}. For the remaining
D-notions, the $s$-way hierarchy of \cite[\S3, Fig.~5 and
App.~A]{BH22} gives, for $\ell \in \{1,3\}$,
$\Advcmtds{\ell}_\TW(\mathcal{A}) \leq
\Advcmtds{4}_\TW(\mathcal{A})$ for the same output distribution. By the
conservative committing-game accounting of \Cref{sec:resources}, the
CMTDs-$\ell$ and CMTDs-4 resource caps charge the same deterministic
decryption checks after syntax, domain, and message-length validation,
independently of their distinctness gates.

For CMTs-$\ell$, the challenger performs $s$ deterministic encryption
checks, which supply the candidate final-root transcripts
(\Cref{lem:root-candidate-seq} builds the candidate sequence for
encryption and decryption checks alike). On a winning execution all
$s$ encryptions are equal, hence of equal length and with one common root
tag $T$. Pairwise $\WiC_1$-distinctness gives pairwise distinct keys,
pairwise $\WiC_3$-distinctness gives pairwise distinct $(K, U, A)$ triples,
and pairwise $\WiC_4$-distinctness gives pairwise distinct full inputs; in
the last case, equal triples and equal ciphertexts would decrypt
deterministically to equal messages, contradicting full-input
distinctness. Thus each CMTs-$\ell$ winner reaches pairwise distinct
opening triples. Opening Triple Separation (\Cref{cor:triple-sep}) makes
the $s$ final-root candidates distinct while their $\tauroot$-bit
projections all equal $T$.
The CMTs-$\ell$ encryption checks are the construction-internal work
counted in $N$ by \Cref{sec:resources}.

By the Lift Application corollary (\Cref{cor:lift-application}) for the
CMTs-$\ell$ game, the derived random oracle adversary $\mathcal{B}$ has
$\qRO(\mathcal{B}) \leq \Nprim$. The Root Tag Hash and Target-Hit Bound
(\Cref{lem:root-tag-hash}, item~1) then bounds its win by
$\RootColl_{\tauroot}(\qRO(\mathcal{B}), s)$ in the random oracle model,
which the lift transfers to the public permutation setting at $\Nprim$;
the encryption checks are construction-internal work counted in $N$, not
in $\Nprim$.
\end{proof}
